\documentclass[11pt]{article}
\newcommand{\tablelang}{en}
\input{../../../paper/tex/preamble}
\usepackage{float}
\newcolumntype{L}{>{\raggedright\arraybackslash}X}
\newcolumntype{R}{>{\raggedleft\arraybackslash}X}
\newcolumntype{P}[1]{>{\raggedright\arraybackslash}p{#1}}
\crefname{appendix}{Appendix}{Appendices}
\setlength{\emergencystretch}{2em}
\hypersetup{pdftitle={When Fitness Scores Mislead: An Empirical Study of LLM-Guided Program Evolution},
  pdfauthor={Oleksandr Pylypenko},
  pdfkeywords={LLM-guided program evolution, OpenEvolve, program search, evaluator validity, benchmark reliability, Gemma}}

\begin{document}
\selectlanguage{english}
\title{When Fitness Scores Mislead: An Empirical Study of LLM-Guided Program Evolution}
\author{Oleksandr Pylypenko\\[2pt]\small Independent Researcher in Artificial Intelligence and Machine Learning\\
\small PhD Student, State University of Trade and Economics, Kyiv, Ukraine\\
\small In cooperation with the Department of Software Engineering and Cybersecurity\\[2pt]
\small ORCID: \href{https://orcid.org/0009-0004-9610-9242}{0009-0004-9610-9242}}
\date{21 November 2026}
\maketitle

\begin{center}\small\textit{Preprint. This manuscript has not been peer reviewed.}\end{center}

\begin{abstract}
Large language model (LLM)-guided program evolution uses an evaluator both to guide search and to select a candidate for use. If the evaluator omits a task requirement, additional search can amplify rather than correct the mismatch. We report two descriptive studies registered before their campaigns, using OpenEvolve 0.4.0 on circle packing and approximate nearest-neighbour search. The Colab L4 study compared Gemma 4 configurations, search modes and budgets; a separate laptop pilot compared evolutionary search and best-of-N with Gemma 4 12B. In three 1,000-call MoE evolutionary runs, the best valid circle packings scored 2.432, 2.204 and 2.258 at the published precision, while the selected verified outcomes were 2.432, 0 and 0. Thus, in two runs, selection discarded a valid improvement found during search. After serial kNN re-timing, the cloud-assisted arm scored above its matched Gemma-only arm in eight of nine pairs, with a median paired difference of 0.776 (exploratory). Among 213 cloud-assisted kNN candidate evaluations ranked above the working seed, 60 (28.2\%) did not constitute verified improvements: four were invalid and 56 were valid but scored no higher than the verified seed. A later, separately frozen six-pair kNN estimation cohort compared a distance-gated evaluator with a timing-matched control: selected programs passed independent checks in 5/6 and 6/6 runs, respectively. This small cohort is descriptive and does not establish an evaluator effect. Overall, the findings are a bounded audit of two tasks and one framework; they do not establish general scaling laws or a causal benefit of cloud assistance.
\end{abstract}

\noindent\textbf{Keywords:} LLM-guided program evolution; OpenEvolve; program search; evaluator validity; benchmark reliability; independent verification.

\medskip
\noindent\textbf{Study independence.} The experiments were conducted by the author using personal equipment and service accounts. The work was not coordinated with or endorsed by Google, OpenEvolve maintainers, or the authors of the cited studies. The author reports the academic status and departmental cooperation stated above.

% =====================================================================
\section{Introduction}\label{sec:intro}

LLM-guided program search couples a proposal mechanism to an evaluator. The model proposes code; an execution-based score guides selection; successful candidates become context for subsequent proposals. Systems such as FunSearch and AlphaEvolve demonstrate the promise of this pattern on problems with executable objectives and high-value algorithmic outcomes~\citep{romera2024funsearch,novikov2025alphaevolve}. Open implementations have extended the design to a broader set of models and tasks~\citep{assumpcao2025codeevolve,sharma2026openevolve}.

The evaluator is part of the search algorithm's effective specification. A candidate can earn a high score while violating a requirement the score does not encode. Repeated sampling then increases exposure to evaluator weaknesses: the candidate selected for shipment is the maximum-scoring one, not necessarily the most correct one. This concern is related to the software-testing oracle problem and to findings that ordinary test suites can miss functional errors in generated programs~\citep{barr2015oracle,liu2023evalplus}.

We examine this issue in two registered descriptive studies of OpenEvolve on two small tasks. The Colab study varies Gemma 4 model configurations, compares evolutionary search (L) with best-of-N sampling (B), includes a cloud-assisted evolutionary arm (LC), and tests longer budgets. The laptop pilot compares L and B with a 12B model. Both tasks have a working evaluator that drives search and an independent post-run re-check. The re-check makes it possible to distinguish a candidate's search score from its validated task outcome.

Our contribution is an empirical audit of search outcomes and evaluator agreement under these conditions. We do not propose a new evolutionary algorithm. Prior work has already evaluated OpenEvolve and related systems on the same $n=26$ circle-packing benchmark; the present study focuses on false improvements, independent verification, model-scale observations, budget effects and the consequences of timing noise~\citep{assumpcao2025codeevolve}. Specifically, we ask:

\begin{itemize}
\item \textbf{RQ1: Model scale and search mode.} How do verified circle-packing scores vary across four Gemma 4 configurations and between L and B at a fixed 100-call budget?
\item \textbf{RQ2: Search budget.} Does a longer search improve the best valid candidate and the score of the candidate selected by the working evaluator?
\item \textbf{RQ3: Cloud assistance.} How do LC and L compare in matched MoE runs on circle packing and kNN?
\item \textbf{RQ4: Evaluator agreement.} How often do candidates ranked above the seed by the working evaluator fail the independent re-check or fail to improve the verified seed score?
\end{itemize}

The original campaign registrations specified descriptive reporting and no significance tests. Their inferential comparisons were added after those campaigns and are explicitly exploratory. The later B2 cohort had a separate protocol frozen before its calls and was designed for estimation only; no confirmatory test was planned or computed. We do not pool its results with the original campaigns.

% =====================================================================
\section{Related work}\label{sec:related}

\paragraph{LLM-guided evolutionary search.} FunSearch combines an LLM proposer with program execution and an evolutionary loop to search over solutions to mathematical and algorithmic problems~\citep{romera2024funsearch}. AlphaEvolve generalizes the approach to scientific and industrial algorithm discovery, using one or more evaluators to guide iterative code changes~\citep{novikov2025alphaevolve}. Open-source frameworks make these methods easier to run with a range of model providers. CodeEvolve includes controlled comparisons against OpenEvolve under matched configurations and budgets, including CirclePackingSquare with $n=26$~\citep{assumpcao2025codeevolve}. That work primarily studies solution quality and algorithmic components. Here we report how often the working fitness score agrees with independent correctness checks in a different model and execution regime.

\paragraph{Evaluation strength and test oracles.} The oracle problem concerns how a test determines whether an output is correct when the expected result is not directly available~\citep{barr2015oracle}. EvalPlus showed that strengthening tests for code-generation benchmarks can expose programs previously counted as correct~\citep{liu2023evalplus}. Google Cloud's AlphaEvolve guidance also documents evaluator exploitation and recommends deterministic checks and held-out validation of winners~\citep{googlecloud2026rewardhacking}. A closely related empirical study, EvoSort, audits correctness gates, ablations and seed sensitivity in LLM-driven sorting-program search, and revises earlier claims after expanding the seed set~\citep{khudhair2026evosort}. These works establish that evaluator reliability and multi-seed auditing are active concerns. Our narrower contribution is a bounded measurement of candidate-level false improvements and selected-run outcomes in OpenEvolve, including a serial timing correction for its kNN task; we do not claim to introduce the general idea of auditing evolutionary search.

\paragraph{Efficiency under correctness constraints.} ECCO evaluates code-generation and editing approaches for improving program efficiency while monitoring functional correctness; its authors also identify timing reliability across hardware as a benchmark challenge~\citep{waghjale2024ecco}. Our kNN task shares that measurement concern, but differs in using an iterative evolutionary search loop and a speed-sensitive working fitness score. The serial re-timing here is a correction for one concrete source of measurement noise, not a substitute for a broader efficiency benchmark.

\paragraph{Task realism.} Repository-level repair benchmarks such as SWE-bench evaluate whether a model-generated patch resolves real GitHub issues in an existing codebase~\citep{jimenez2024swebench}. The present experiments use small bounded tasks with explicit, executable evaluators. This narrower design makes candidate-level re-checks tractable but leaves external validity to future work on real repositories.

% =====================================================================
\section{Study design and measures}\label{sec:design}

\subsection{Search conditions}

All campaigns used OpenEvolve 0.4.0 with the registered task seeds and configuration~\citep{sharma2026openevolve}. In \textbf{L}, Gemma generates mutations within the evolutionary population and archive. In \textbf{B}, the same model samples independent candidates from the starting program without evolutionary history; the working evaluator ranks the candidates. In \textbf{LC}, the evolutionary loop is retained and approximately 10\% of mutation calls are routed to Gemini 4 Argon, with the remaining calls made by Gemma. The call budget is a count of logged LLM calls, not a common unit of compute, latency or energy.

\subsection{Tasks and working evaluators}

\textbf{Circle packing.} The task places 26 circles inside a unit square without overlap and seeks a large sum of radii. The working evaluator uses a normalized sum of radii divided by 2.635, the reference value used by the benchmark, and assigns zero to candidates it rejects. We do not treat 2.635 as a mathematically proven upper bound. The independent re-check requires finite positive radii, containment, pairwise non-overlap at a tighter tolerance, deterministic output across two sandbox executions and compliance with static restrictions. The paper reports the re-check's raw sum of radii. The seed itself is invalid because two radii are zero, so its verified score is zero.

\textbf{Approximate k-nearest neighbours.} The working evaluator combines recall and relative speed against exact search, but does not check that the returned squared distances correspond to the returned neighbour indices. The independent re-check uses fresh challenge inputs, checks exact-copy queries and distance correctness, and evaluates speed and recall on separate holdout sets. Consequently, a candidate can achieve high working fitness through a precision shortcut while failing the task's distance requirement.

\subsection{Independent endpoints}

For run $r$, let $s_w(c)$ be the working evaluator's score for candidate $c$, and let $s_v(c)$ be the independent re-check score when the candidate passes validity. The selected candidate is $c_r^*=\arg\max_c s_w(c)$. The primary endpoint is its re-check score; if it fails validity, the run receives the seed's verified score. A \emph{false improvement} is a candidate evaluation with $s_w(c)>s_w(c_0)$ that is invalid on re-check or has $s_v(c)\leq s_v(c_0)$, where $c_0$ is the seed. The pooled false-improvement denominator counts candidate evaluations, including repeat evaluations of the same code hash, rather than unique programs.

\begin{figure}[tb]
\centering\includegraphics[width=\linewidth]{fig1_pipeline_en.pdf}
\caption{Search and evaluation pipeline. L, LC and B generate candidate sets. The working evaluator guides selection; the independent re-check determines the reported verified endpoint. The holdout is not exposed to candidate generation or working-score selection.}
\label{fig:pipeline}
\end{figure}

\subsection{Campaigns and protocol deviations}

The initial Colab registration, made on 30 October before any study run, planned 156 runs and described descriptive analysis without significance tests. The plan was revised during the campaign: on 5 November it fell to 145 after removing fourteen 31B slots and adding three exploratory long-budget best-of-N circle-packing runs; on 6 November four 31B circle-packing slots were restored, bringing it to 149. On 7 November a documented amendment expanded the plan to 179 by adding 20 LC runs and 10 unpaired L controls, for 30 LC runs total. At that decision, 139/149 original-plan runs were complete and LC runs were already underway; prior L/B outcomes and early LC operational logs had been inspected. The amendment's size was a resource allocation, not a power calculation. The final inventory completed 167/179 runs: 18 LC runs completed, ten were not started and two remained partial after quota limits. The 100-call size ladder used ten seeds per arm for E4B, 12B and 26B A4B on circle packing and kNN, and five seeds per arm for 31B on circle packing. It also included three 1,000-call L circle-packing runs, three 500-call L kNN runs and three 1,000-call B circle-packing runs; the long B arm was exploratory. LC has nine complete matched pairs per task: five in the original cohort and four in the extension cohort. The ten new L controls are unpaired and are excluded from the matched comparison and size ladder.

The separate laptop pilot was registered on 29 October for 18 runs. All 12 planned L/B runs were completed (three seeds per task and arm); the six-run LC condition did not enter the analysis after quota exhaustion, leaving one partial run. The laptop used an RTX 4060 Laptop GPU with 8 GB of VRAM and 31 GB of RAM. Colab used one L4 GPU with 24 GB VRAM and 12 vCPUs. Other deviations and protocol details are recorded in the two preregistrations, the dated LC amendment and the deviation logs~\citep{pylypenko2026artifacts}.

The Colab and laptop 12B runs are not a hardware-controlled comparison. They used different model builds, Ollama versions and candidate-execution environments. In the Colab 12B B condition, the build repeatedly generated an inapplicable search-and-replace block; no candidate was produced in the circle-packing runs. This known failure must be considered when interpreting the 12B L--B contrast.

\subsection{Subsequent B2 estimation cohort}

After the original campaigns, a separate six-pair estimation cohort compared two kNN evaluators under a protocol frozen in Git commit \code{3e603e426930215838ce9a2aa475959bb241ff87} before model calls. Each pair used the same task seed and seed program, local Gemma 4 12B (model digest \code{6114515d63c1}), prompt, sandbox, and 20-call budget per arm. V1-matched used the V2 timing harness but omitted its pre-fitness distance-correctness gate; V2 applied that gate before timing and scoring. Pair order was fixed in the configuration. The run-level endpoint was the official OpenEvolve-selected program, identified by its recorded hash, not the candidate-log maximum.

The frozen protocol was \code{paper2-b2-estimation-2026-11-20-v1} (SHA-256 \code{8f990fd64458c8060c79d71655a9712eb9f9f5963a653482e96452ff1c762719}). The same independently implemented checker evaluated both arms using fixed holdout seeds 9001--9003. These seeds were source-visible and were not supplied to the candidate prompt or working evaluator; the checks are reproducible but not a secret or sealed holdout. The cohort was explicitly estimation/feasibility work, not a confirmatory test. It was analyzed separately from the original campaigns.

\subsection{kNN timing correction}

The first kNN re-check pass ran several Docker sandboxes concurrently. Under that load the identical seed code received scores from 0.02 to 1.0 because the score includes a wall-clock speed component. We therefore serially re-timed the seed and all valid above-seed candidates on an otherwise idle laptop, one sandbox at a time. The seed code hash has eight serial observations: five designated calibration rows and three additional observations from runs in which the same code was selected as the shipped program. The main endpoint pools all eight (median 0.2235; range 0.1606--0.2438); the five calibration rows alone have median 0.238 (range 0.1919--0.2438). Pilot L false improvements change from 36/42 under the calibration-only sensitivity to 33/42 under the pooled median. LC kNN remains 60/213, and LC is higher than L in 8/9 paired runs under both definitions.

The serial pass covered 214 distinct valid above-seed hashes from Colab and 14 additional pilot hashes. It did not retime approximately 1,900 candidates below the working seed score. Accordingly, the registered all-candidate best-verified kNN endpoint is not reported. Among 71 hashes with repeated serial measurements, the median within-hash score range was 0.018; 16 exceeded 0.05 and the largest was 0.490. The re-timing improves comparability for the measured candidates; residual variation remains near the score ceiling.

\subsection{Statistical analysis}

The two original campaign registrations specify descriptive analysis only; the B2 protocol separately prohibits confirmatory inference. We report post hoc two-sided comparisons for the original campaigns as exploratory. Circle-packing outcomes across the four L-arm model groups are summarized with a Kruskal--Wallis test and pairwise Mann--Whitney tests with Holm adjustment. L--B comparisons are also reported per task and model, with Holm adjustment across the nine available comparisons. LC--L analyses use matched seeds, paired sign tests, Wilcoxon signed-rank tests and percentile bootstrap intervals from 10,000 resamples of the paired differences. The B2 cohort has six pairs and is summarized with counts and descriptive scores only. These analyses do not establish equivalence, causal effects or general performance outside these tasks.

% =====================================================================
\section{Results}\label{sec:results}

\subsection{Circle-packing scores across model configurations}

Table~\ref{tab:scale} reports the median shipped verified score at 100 calls. The L-arm medians were 1.63 for E4B and 12B, 2.17 for 26B A4B MoE, and 2.25 for 31B. The overall group difference was detectable in an exploratory Kruskal--Wallis test ($p=0.0155$). None of the three pairwise MoE comparisons passed Holm correction: adjusted $p=0.063$ versus E4B, $p=0.173$ versus 12B, and $p=0.426$ versus 31B. In particular, the MoE--31B result does not demonstrate equivalence.

\begin{table}[tb]
\centering\small
\caption{Circle-packing shipped verified scores in the 100-call size ladder. Scores are raw sums of radii; the independent check determines validity. The L--B p-values are exploratory and Holm-adjusted over nine comparisons across tasks and models.}
\label{tab:scale}
\begin{tabularx}{\linewidth}{@{}lrrrrr@{}}
\toprule
Model & L median & B median & $n$ per arm & L--B $p$ & Holm $p$ \\
\midrule
Gemma 4 E4B & 1.628 & 0.851 & 10 & 0.491 & 1.000 \\
Gemma 4 12B & 1.633 & 0.000 & 10 & 0.0022 & 0.0199 \\
Gemma 4 26B A4B MoE & 2.170 & 1.911 & 10 & 0.0373 & 0.299 \\
Gemma 4 31B & 2.249 & 2.400 & 5 & 0.151 & 1.000 \\
\bottomrule
\end{tabularx}
\end{table}

The 12B circle-packing contrast is not evidence that L generally outperforms B: the Colab B build failed to apply its diffs, leaving all ten runs at the invalid seed score. For 31B, B's median exceeded L's in this five-seed sample. The other medians are descriptive; these results do not support a universal L-over-B conclusion. The Colab 100-call run used about 0.9 compute units for the MoE and 3.9 for 31B. The 31B model was served one run at a time while other models used two parallel runs, so this is an approximate platform-specific cost comparison, not a measurement of energy or equal-throughput efficiency. E4B used slightly fewer compute units per median score point (about 0.40 versus 0.43 for the MoE), while its absolute median score was lower.

\begin{figure}[tb]
\centering\includegraphics[width=0.92\linewidth]{fig2_circle_scale_en.pdf}
\caption{Per-run circle-packing scores for the 100-call L and B size ladder. Points are independently checked run outcomes; dark horizontal marks are medians. The 12B B scores are zero because the Colab build did not apply a candidate diff.}
\label{fig:scale}
\end{figure}

\subsection{Longer search and selected-program validity}

The three 1,000-call MoE L runs found valid candidates with scores 2.432, 2.204 and 2.258 as reported to three decimals. The selected verified outcomes were 2.4317, 0 and 0; the two zero endpoints are the registered fallback to the invalid seed after selection of an invalid candidate. The approximate selection regrets at the source precision are therefore 0.000, 2.204 and 2.258 (median 2.204). In the exploratory 1,000-call B control, the shipped candidates scored 1.936, 2.072 and 1.924 and were valid; the report records the best-valid candidates as the same outcomes at its precision. The L-long false-improvement share was 607/1,064 candidate evaluations (57\%); B-long was 358/2,089 (17\%). This three-seed comparison shows that search found strong valid candidates that the selection rule failed to retain in two runs. It does not establish that long evolutionary search is generally worse than best-of-N. Figure~\ref{fig:selection-audit} shows the gap between best-valid and selected outcomes and the paired kNN results.

For the 500-call MoE L kNN runs, the working evaluator ranked no candidate above the seed, so all three runs shipped the seed. The registered all-candidate best-verified kNN endpoint cannot be reconstructed from the serial pass because candidates below the working seed were not all re-timed.

\subsection{Cloud-assisted MoE runs}

Table~\ref{tab:lc} compares nine matched LC--L runs for each task. In circle packing, LC scored higher in seven pairs, tied in one and scored lower in one. The median scores were 2.630 for LC and 2.192 for L; the exploratory Wilcoxon test was not significant ($p=0.148$). In kNN, after serial re-timing, LC scored higher in eight of nine pairs. Its median was 1.000 versus 0.223 for L, and the median paired difference was 0.776 (95\% bootstrap interval 0.192--0.777). The unadjusted exploratory sign test was $p=0.039$ and Wilcoxon test $p=0.0078$. These comparisons were not pre-registered and do not identify a general causal effect of cloud assistance. The 18 completed LC runs used 198 Argon calls and 4.64 million tokens. All nine selected LC programs passed the independent distance and exact-copy checks; their median holdout recall ranged from 0.99985 to 1.000. The median of each program's three serial maximum holdout speed ratios was below 1.0 in eight runs and 1.255 in the remaining run. These are relative ratios, not absolute latency measurements. Across the ten unique selected code hashes (the shared seed plus nine LC programs), pairwise Kendall's $\tau_b$ for the three ordinal repeat indices ranged from 0.762 to 0.810; top-five overlap was four or five programs. No LC program crossed the pooled seed-score threshold across its three repeats. This descriptive stability check is limited to selected programs; repeat indices are not synchronized across programs.

Descriptively, the five initially registered pairs (seeds 101--505) had LC higher/equal/lower counts of 3/1/1 for circle packing (median L 2.148; LC 2.626) and 5/0/0 for kNN (L 0.223; LC 1.000). In the four later matched pairs (606--909), the corresponding counts were 4/0/0 and 3/0/1, with medians of L 2.215 versus LC 2.634 for circle packing and L 0.223 versus LC 0.708 for kNN. These subgroup summaries are descriptive only; the nine-pair comparison remains the primary presentation of this exploratory condition.

\begin{table}[tb]
\centering\small
\caption{Matched MoE LC--L outcomes. Each comparison contains nine seeds. All inferential tests are post hoc and unadjusted.}
\label{tab:lc}
\begin{tabularx}{\linewidth}{@{}lrrrrL@{}}
\toprule
Task & L median & LC median & LC higher / equal / lower & Wilcoxon $p$ & Note \\
\midrule
Circle packing & 2.192 & 2.630 & 7 / 1 / 1 & 0.148 & No reliable direction in this small sample \\
kNN, serial timing & 0.223 & 1.000 & 8 / 0 / 1 & 0.0078 & Exploratory \\
\bottomrule
\end{tabularx}
\end{table}

\begin{figure}[tb]
\centering\includegraphics[width=0.95\linewidth]{fig3_lc_pairs_en.pdf}
\caption{Paired MoE outcomes for Gemma-only L and cloud-assisted LC. A line joins the same seed. The reported p-values are exploratory, unadjusted signed-rank tests.}
\label{fig:lc}
\end{figure}

\begin{figure}[tb]
\centering\includegraphics[width=0.98\linewidth]{fig4_selection_audit_en.pdf}
\caption{Selection audit. Left: published best-valid scores and selected verified endpoints for the three 1,000-call L runs and exploratory B controls. Best-valid values are available only to three decimals in the public summary. Right: selected kNN scores for nine matched L--LC pairs, recomputed from serial re-checks. The 8/9 direction is exploratory.}
\label{fig:selection-audit}
\end{figure}

\subsection{Laptop pilot}

On the laptop, L shipped a valid circle packing in all three runs (median 1.773; range 1.596--2.120). B shipped an invalid program in two of three runs and had median primary score zero (range 0--1.793). With only three seeds, this is descriptive; the exploratory Mann--Whitney test gives $p=0.376$. In all six pilot kNN runs, the top working-score candidate failed the independent check, so L and B both fell back to the seed score of 0.223. The failures returned high-recall neighbours but used float32 arithmetic that violated the specified true-squared-distance tolerance. The working evaluator did not check that requirement.

The laptop and Colab 12B results cannot isolate a hardware effect. In addition to different execution environments, the Colab 12B build differed from the laptop build. The Colab 12B B failure described above is a concrete example of how a model build can interact with a diff-based search loop.

\subsection{B2 selected-program estimation results}

All six B2 pairs completed, yielding 12 runs and 240 logged model calls. There were 207 calls marked \texttt{ok} and 33 token-limit truncations; every truncated response reached the 2,048-token completion cap. The pinned independent re-check verified all 12 selected artifact hashes against their run manifests before evaluation. Table~\ref{tab:b2} reports the per-pair independent score and validity of the program selected by each evaluator. Scores for invalid selections are set to zero under the frozen analysis rule.

\begin{table}[H]
\centering\small
\caption{B2 selected-program outcomes in the separate six-pair estimation cohort. Scores are independent holdout scores; an invalid selection is scored zero.}
\label{tab:b2}
\begin{tabularx}{\linewidth}{@{}lrrll@{}}
\toprule
Pair & V1-matched score & V2 score & V1-matched valid & V2 valid \\
\midrule
P1 & 0.2107 & 0.2491 & Yes & Yes \\
P2 & 0.2471 & 0.2141 & Yes & Yes \\
P3 & 0.2310 & 1.0000 & Yes & Yes \\
P4 & 0.2043 & 0.2588 & Yes & Yes \\
P5 & 0.2241 & 0.0000 & Yes & No \\
P6 & 0.2346 & 0.2542 & Yes & Yes \\
\midrule
Valid selections & 6/6 & 5/6 & 100\% & 83.3\% \\
\bottomrule
\end{tabularx}
\end{table}

The paired validity table contains five pairs in which both selections passed, one in which V1-matched passed and V2 failed, and none in which only V2 passed or both failed. The descriptive difference in invalid-selection rates (V2 minus V1-matched) is +16.7 percentage points. Mean selected score, assigning zero to invalid selections, was 0.2253 for V1-matched and 0.3294 for V2. The six paired score differences (V2 minus V1-matched, with invalid selections assigned zero) were +0.0385, $-0.0330$, +0.7690, +0.0545, $-0.2241$ and +0.0195 for P1--P6. Their median was +0.0290 (range $-0.2241$ to +0.7690); this descriptive distribution is sensitive to the P3 difference of +0.7690 and is not evidence of an evaluator effect. The only invalid selection was P5 V2: its working evaluator had marked it valid, but the independent checker reported 140 fresh-query failures because returned squared distances were not derived from the input vectors. Holdout recall for the selected program was 1.0, illustrating why recall alone did not establish contract compliance.

No confirmatory p-value was computed. With six pairs, one discordant pair, one task, one local model configuration, and a source-visible holdout, this cohort neither demonstrates that V2 reduces invalid selection nor supports equivalence or generalization. Account-level Colab usage was pooled with other assignments, so its 35.56-CU balance change from launch cannot be attributed to B2 alone.

\subsection{False improvements under independent re-check}

Table~\ref{tab:false} shows that every evaluated group with above-seed candidates had at least some false improvements, with rates depending on task, arm and budget. In the MoE LC kNN condition, 60 of 213 above-seed candidate evaluations were false improvements (28.2\%): four were invalid and 56 passed validity but had a serially verified score no greater than the pooled seed score. This classification does not equate score with latency alone: the score combines recall and a relative-time penalty. In the laptop pilot, 33/42 L kNN evaluations were false (79\%), while B had five above-seed evaluations and all five were false. These denominators count candidate evaluations, not unique programs; the pilot B rate is especially unstable because $n=5$.

\begin{table}[tb]
\centering\small
\caption{False improvements among candidate evaluations ranked above the seed by the working evaluator. Entries are false / above-seed evaluations (rate). A dash means no candidate was ranked above the seed.}
\label{tab:false}
\begin{tabularx}{\linewidth}{@{}lrr@{}}
\toprule
Study condition & Circle packing & kNN \\
\midrule
E4B L & 239 / 505 (47\%) & 1 / 2 (50\%) \\
E4B B & 213 / 408 (52\%) & 0 / 1 (0\%) \\
12B L, Colab & 103 / 423 (24\%) & 6 / 9 (67\%) \\
12B B, Colab & --- & 1 / 5 (20\%) \\
26B MoE L & 129 / 421 (31\%) & 1 / 1 (100\%) \\
26B MoE B & 117 / 710 (17\%) & --- \\
26B MoE LC & 140 / 486 (29\%) & 60 / 213 (28\%) \\
31B L & 56 / 404 (14\%) & not run \\
31B B & 90 / 336 (27\%) & not run \\
MoE L, long budget & 607 / 1,064 (57\%) & --- \\
MoE B, long budget & 358 / 2,089 (17\%) & not run \\
12B laptop L & 21 / 140 (15\%) & 33 / 42 (79\%) \\
12B laptop B & 22 / 86 (26\%) & 5 / 5 (100\%) \\
\bottomrule
\end{tabularx}
\end{table}

Circle-packing false improvements often exploited differences between the working score and the stricter validity criteria, including zero-radius circles. In kNN, the working evaluator's omission of returned-distance correctness allowed fast approximate candidates to be rewarded. The false-improvement share rose to 57\% in the long MoE L circle-packing runs. This association is consistent with more search exposing the evaluator's gaps, but the comparison is exploratory and the longer runs also differed in budget and candidate distribution.

% =====================================================================
\section{Discussion}\label{sec:discussion}

\paragraph{Fitness is an operational objective, not a correctness proof.} Across both tasks, candidates can be ranked above the seed by the working evaluator and still fail an independent check. The kNN case is particularly clear: recall and speed are insufficient when the API contract also requires accurate squared distances. The separate B2 estimation cohort shows that a pre-fitness correctness gate is implementable, but selected-program validity was 5/6 under V2 and 6/6 under the timing-matched control. The single V2 failure illustrates a remaining gap; the small cohort does not establish that the gate improves or worsens selection reliability.

\paragraph{Search budget can amplify evaluator defects.} The long MoE L runs found strong valid circle packings, but the working evaluator selected invalid candidates in two of three runs. More calls therefore improved the best-valid candidate pool without reliably improving the selected output. This distinction suggests that future studies should report both best-valid and selected-verified outcomes, and should audit the exact program that the search process would ship.

\paragraph{Model scale and compute.} The observed L-arm medians increase between the smaller configurations and the MoE/31B runs, but the pairwise size comparisons do not survive the reported correction. The 31B sample has five runs; the MoE-versus-31B test neither establishes a difference nor equivalence. The compute-unit comparison also reflects different serving concurrency. These data motivate larger matched studies with consistent serving, model-build digests and energy or cost accounting.

\paragraph{Cloud assistance.} The MoE LC arm had higher kNN verified scores than L in eight of nine pairs after the timing correction. The post-hoc nature, small number of pairs, partial completion of the planned LC runs and extra cloud calls limit the strength of this finding. Circle-packing results were directionally favorable to LC but inconclusive. A follow-up should pre-register the comparison, repeat it across independent tasks and models, and include cloud cost in a common unit.

\paragraph{What this paper adds.} Prior work has shown that LLM-guided evolutionary search can discover high-performing programs, including on the circle-packing family~\citep{romera2024funsearch,novikov2025alphaevolve,assumpcao2025codeevolve}. Our evidence addresses a complementary operational question: how often does a search evaluator's ranking survive independent verification, and how does that answer change after timing is measured serially? The later B2 cohort adds a small, separately reported estimate for selected-program validity under a distance-correctness gate. Together, these are bounded audits of one framework and model family, not a general benchmark of coding agents or evidence of a general gate effect.

% =====================================================================
\section{Limitations and threats to validity}\label{sec:limitations}

\begin{itemize}
\item \textbf{Task scope.} The study uses two small, constructed tasks. Circle packing is an established benchmark, not a novel problem; approximate kNN is a bespoke task. Results do not establish performance on real repository defects or scientific code.
\item \textbf{B2 estimation limits.} The later evaluator comparison contains six pairs on one kNN task and one local model digest. It was designed for estimation, not power to detect a treatment effect; the source-visible fixed holdout is not sealed. The observed 6/6 versus 5/6 validity counts are imprecise and do not establish efficacy or equivalence. The aggregate Colab meter was shared with other assignments, preventing exact B2-only compute-unit attribution.
\item \textbf{Model scope.} The scale ladder uses several configurations within Gemma 4. The cloud-assisted comparison uses one Argon setting. Model quality, architecture and deployment environment vary together; the experiments do not identify a general scaling law.
\item \textbf{Incomplete campaigns.} The Colab study completed 167/179 amended-plan runs. The laptop pilot completed all 12 planned L/B runs; its separate six-run LC condition was excluded after quota exhaustion, with one partial LC run. The pilot has only three seeds per task and arm.
\item \textbf{Exploratory inference.} Both registrations specified descriptive analysis. The tests reported here were added after results were known; some p-values are unadjusted and the Holm corrections apply only to the stated comparison families. Small sample sizes limit precision.
\item \textbf{Model-build and environment confounding.} The Colab 12B B arm suffered a build-specific diff-application failure. Cross-device 12B comparisons additionally vary model digest, Ollama version and execution environment. The 31B cost comparison used different concurrency.
\item \textbf{Evaluator and holdout limitations.} The independent re-check is stronger than the working evaluator but is not a proof of correctness. It covers specific constraints and finite holdout inputs. Search may exploit requirements neither evaluator captures.
\item \textbf{Timing sensitivity.} kNN fitness includes wall-clock speed. Serial re-timing was performed for the seed and all valid above-seed candidates, not every generated candidate. Approximately 1,900 below-seed candidates remain untimed, and repeated valid programs still show residual score variation.
\item \textbf{Data access.} Preregistrations, deviations, run-level analysis, serial retiming and code are public. Complete model-call logs and all raw run directories are archived separately and are not included in the public Git snapshot.
\end{itemize}

% =====================================================================
\section{Data and code availability}\label{sec:availability}

The original campaign preregistrations, dated amendments, deviation logs, run-level analysis, serial kNN re-check files and code are available in the immutable repository snapshot at commit \code{57a55d3d23d6c7c3dd7be3abfc1a0d9c758fc6b5}~\citep{pylypenko2026artifacts}. The separate B2 estimation protocol was frozen before calls in source commit \code{3e603e426930215838ce9a2aa475959bb241ff87}; its configuration, final report, selected-program re-checks and descriptive summary are available at \code{experiments/openevolve/configs/b2\_estimation\_2026-11-20.json} and \code{docs/evaluator\_v2/}. The per-run B2 re-check and summary are committed as \code{B2\_SELECTED\_RECHECK\_2026-11-21.json} and \code{B2\_ESTIMATION\_SUMMARY\_2026-11-21.json}. Full raw run folders and model-call traces are retained in a separate GCS archive and are not included in the Git snapshot. The original analysis and figures are rebuilt by \code{bash scripts/paper\_ii/build.sh}. No archival DOI is claimed.

% =====================================================================
\section{Conclusion}\label{sec:conclusion}

Across the two original descriptive studies and the separate B2 estimation cohort, independent checks exposed gaps between working fitness and verified behavior. In two of three 1,000-call MoE L runs, the search found a valid circle packing that the selection rule did not retain. The post hoc serial kNN analysis found a higher verified score in eight of nine LC--L pairs and 60/213 false-improvement candidate evaluations above the seed. In B2, one of six V2-selected programs failed the distance contract, while all six timing-matched V1 selections passed; this small descriptive comparison does not establish a gate effect. Incomplete campaigns, narrow task scope, timing sensitivity, a build-specific failure and small estimation samples limit broader claims. Future work should compare evaluator designs on independent tasks with adequately planned samples and report raw traces, verified outputs and end-to-end costs.

\section*{Research transparency}
\textbf{Funding and competing interests.} No external funding is reported. The study used the author's hardware and service accounts. The author declares no competing financial interest related to OpenEvolve or the cited systems.

\textbf{Institutional relationship.} The author reports PhD-student status at the State University of Trade and Economics and cooperation with its Department of Software Engineering and Cybersecurity. This statement describes the author's reported status and cooperation; it does not imply that the department ran or funded the experiments.

\textbf{Use of AI tools.} AI assistants supported analysis-script development, manuscript drafting and Ukrainian translation. They did not conduct the model campaigns. The author is responsible for verifying the analyses, citations, translations and final claims. This draft has not received independent human peer review.

\bibliographystyle{plainnat}
\bibliography{../references}

\clearpage
\appendix
\section{Registration and deviations}\label{app:registration}

The laptop pilot was registered on 29 October under tag \code{oe-pilot-prereg}; the Colab scale study was registered on 30 October under tag \code{colab-scale-prereg}. The initial Colab plan was 156 runs. On 5 November, after interim results and budget had been reviewed, fourteen 31B slots were removed and three exploratory long-budget best-of-N circle-packing runs were added (145 total). On 6 November four 31B circle-packing slots were restored (149 total). On 7 November, after 139/149 runs were complete and while initial LC runs were in progress, the owner amended the ongoing plan by adding 20 LC runs and 10 unpaired L controls (179 planned; 30 LC runs total). Earlier L/B results and LC operational logs had already been inspected. The amended LC size was a resource allocation, not a power calculation. The final inventory contains 167 completed runs, ten unstarted LC runs and two quota-paused partial LC runs. The pilot completed all 12 planned L/B runs; the six-run LC condition was excluded after quota exhaustion. The separate B2 estimation protocol was frozen in Git commit \code{3e603e426930215838ce9a2aa475959bb241ff87} before its first calls; it specified six matched kNN pairs and descriptive analysis only. The B2 cohort completed all 12 arms, and no sample-size expansion or confirmatory analysis was performed. Other material deviations include the Colab 12B B diff-application failure, serial re-timing after concurrent measurements of the seed ranged from 0.02 to 1.0, and post hoc inferential tests for the original campaigns. Full dated records are in the preregistrations, protocol and deviation logs.

\section{Evidence map}\label{app:evidence}

\begin{minipage}{\linewidth}\small
\captionof{table}{Primary manuscript claims and their committed evidence sources.}
\label{tab:evidence}
\begin{tabularx}{\linewidth}{@{}P{0.33\linewidth}L@{}}
\toprule
Claim or analysis & Evidence \\
\midrule
Colab run status and scale results & \path{docs/COLAB_SCALE_RESULTS.md}; \path{docs/colab_scale/FINAL_RUNS_2026-11-11.json} \\
Pilot results & \path{docs/OPENEVOLVE_PILOT_RESULTS.md}; \path{docs/openevolve_pilot/summary_2026-11-19.json} \\
Registered design and departures & \path{docs/colab_scale/PREREGISTRATION.md}; \path{docs/colab_scale/DEVIATIONS.md}; \path{docs/colab_scale/LC_EXTENSION_2026-11-07.md} \\
Serial kNN correction & \path{docs/colab_scale/results/knn_serial_retiming_2026-11-19.jsonl}; \path{knn_serial_above_seed_2026-11-19.jsonl} \\
Selection audit & \path{docs/paper_ii/results/revision_metrics.json}; \path{scripts/paper_ii/revision_metrics.py} \\
Exploratory statistics & \path{docs/colab_scale/results/exploratory_tests_2026-11-19.json}; \path{scripts/analysis/oe_exploratory_tests.py} \\
Article figures & \path{scripts/paper_ii/make_figures.py}; \path{fig4_selection_audit_en.pdf} \\
B2 frozen protocol and configuration & \path{experiments/openevolve/configs/b2_estimation_2026-11-20.json}; \path{docs/evaluator_v2/B2_ESTIMATION_DRAFT.md} \\
B2 selected-program re-check and descriptive summary & \path{docs/evaluator_v2/B2_SELECTED_RECHECK_2026-11-21.json}; \path{docs/evaluator_v2/B2_ESTIMATION_SUMMARY_2026-11-21.json} \\
\bottomrule
\end{tabularx}
\end{minipage}

\end{document}
